\section{Cohort, instrument and denominators}
The quantitative reference cohort is the historical 423 records. In the saved coded accounting, the key filter removed no records (490 remained); specified consent codes left 486, gender codes left 478, and remittance codes left 423. The original manual test and non-official exclusions also remain in force. Their precise position relative to that export is not established by the coded counts. This is not a response rate, an exhaustive eligibility check or a verified count of distinct households. Referral and household links were not systematically recorded.

The five submitted form versions are 2512090157, 2512091719, 2512111044, 2512121811 and 2602120658. Unresolved versions stay unresolved. Early residence code 4 means Other; later code 4 means Soacha. Meta is an additional later option and does not identify a municipality. Use recorded dates and distinct recall windows, not inferred interview sites, exposure or migration tenure.

Keep route (open, closed or unresolved) distinct from response state: ordinary valid answer, structural skip, unavailable item or option, applicable nonresponse, item-specific sentinel, or unresolved state. A recorded-answer base may include an offered sentinel; numeric summaries use scorable answers. Multiple-response shares use the relevant valid-parent base and versions offering the option. An unoffered option, an undefined share, a withheld cell and an observed zero are different. A diagnostic flag does not alter the cohort.

\section{Outcomes and focal fields}
\begin{description}
\item[Remittance disposition (q6\_13).] The household screen q5 asked about international remittances received by the respondent or someone in her household over the preceding 12 months. On the eligible route, q6\_13 was asked independently of q6, the principal-recipient or manager question. The five choices collapse to reported all cash (1) versus any reported digital retention or use (2--5). Codes 2 and 4 overlap; the item does not measure pesos, digital spending alone, the latest receipt channel, or a withdrawal following digital receipt. The base is the household-screened referral cohort with valid answers rather than a sample verified to receive remittances personally.
\item[Domestic digital participation (q6\_14).] Valid No/Yes for sending or receiving within 60 days; refusal and nonresponse are unscored. The q6\_15 follow-up is conditional on Yes. The five operation bins use the valid-bin denominator, whereas response-state accounting uses all affirmative users. Uncertainty (98) is not a bin; the grouped answers do not yield exact counts.
\item[Remittance context.] q6\_2 records 12-month receipt frequency; combined code 99 is not simple nonreceipt. q6\_4 records the latest receipt channel with Other unspecified. Each has its own recorded-answer base. The inbound items q6\_6 and q6\_9 ask separately about time from sending to usable funds and the clarity of any exchange rate applied to the most recent remittance. Domestic fee visibility (q6\_18) concerns the last digital operation in a 90-day window and is asked only on the affirmative 60-day-use branch (q6\_14). These are different items; neither routing nor wording supplies an outcome distribution.
\item[Phone, internet and data access (IAT3).] Fixed mean of q3\_1, q3\_5 and q3\_8. Phone-use codes 1--3 score $1,0.5,0$; internet-frequency codes 1--5 score $1,0.75,0.5,0.25,0$. Codes 1--3 of q3\_8 score $0,0.5,1$ and describe whether running out of data or balance prevented use of financial services or communication in the preceding month. That item was offered only to regular smartphone or basic-phone users (q3\_1 codes 1 or 2). The phone component's code 3 can score zero but cannot enter a complete score through this route. Missing and other unscorable components leave the full score unscored.
\item[Reported transfer ability and fraud-avoidance confidence (IADT2).] Equal-weight average of reported phone-transfer ability (q3\_16, also offered only to regular smartphone or basic-phone users) and reported confidence in avoiding online fraud or suspicious messages (q3\_17). IADT3 adds a third equally weighted item: reported ability to install an app (q3\_15). That item is asked only on the personal smartphone route (q3\_1, code 1). Each component maps an ordinary eligible response code \(c=1,\ldots,5\) to \((c-1)/4\). Both scores require all their designated components, and neither is an observed ability test.
\item[Reported onboarding conditions (IAFF2).] Mean of reporting own-name address proof described by the respondent as acceptable to financial institutions (q4\_5) and own-name registration of the number normally used for digital accounts or wallets (q4\_6). Possible values are 0, 0.5 and 1. A provider's decision, legal eligibility and account ownership are not verified.
\end{description}
Experiences, safety, conditional recourse, household decisions, training, support and barriers remain separate item families. A pair-complete comparison uses records with the two measures being compared; a common-complete comparison fixes one base across the specified set. Shared items are not independent corroboration. The supplement reports only the comparisons identified in its tables and figures.

\section{Conditional models and uncertainty}
The five working logit models share IAT3, IADT2, IAFF2, age, grouped education (completed secondary as reference) and grouped recorded residence (Bogotá as reference). R0 and D0 use the remittance and domestic outcomes, respectively. R1 and D1 add occupation (employee or worker as reference) and the recorded arrival-year field (q2\_1), not tenure. The question asks about first arrival, but for a woman who entered more than once the enumerator instruction was to record the most recent arrival to stay continuously. The residual residence group combines Other, Soacha and Meta where recorded; it is not a city or cluster. R0C fits the remittance core on the intersection of the two intended core samples, which equals the domestic core sample.

The core sample requires a valid outcome, all three focal measures, resolved residence, an ordinary recorded age response with an integer value from 18 to 100, and a valid education code (1--10). The extensions additionally require valid occupation (1--9) and recorded arrival year (1950--2026). These are model input checks, not new exclusions from the admitted cohort. The five fitted samples match their intended rows. R1 uses the same rows as R0, and D1 uses the same rows as D0. R0C keeps the remittance outcome and core specification fixed and changes only the analysis sample relative to R0.

The contrasts average predictions at the two fixed score values over each entire fitted sample, leaving other covariates as observed. IAT3 and IADT2 compare 0.50 with 0.75; IAFF2 compares 0.50 (either one of the two conditions, pooled) with 1.00 (both conditions). The differences are in percentage points. Using one slope on the IAFF2 score gives phone-only and proof-only the same prediction at fixed covariates. The education-by-residence screen distinguishes lack of exact endpoint observations from endpoints outside the observed range. Its denominator counts strata, not people; the screen neither establishes overlap on all covariates nor trims the averages.

Pointwise 95\% delta-method intervals use the working robust covariance assuming independent records. The six separate coefficient-null Wald tests form one fixed Holm family. No causal, design-weighted, referral-clustered or simultaneous inference is implied.
